\documentclass[11pt]{article}
\usepackage{amsmath,amssymb,amsthm}
\newcommand{\Cov}{\operatorname{Cov}}
\newcommand{\Var}{\operatorname{Var}}
\title{The actual mixed-transition covariance interface: retained histories and old currents}\author{}\date{}
\begin{document}\maketitle
\paragraph{Scope and sources.}
The sources are the Wilson-source estimates \cite{W} and \cite{ND}.
The target is (MT36) of \cite[p.~692]{W}, including its lower/contact companions, for a
completed stable birth $R_j$ with its dimension-six projection removed:
\[
 \|\mathcal P_6\mathcal L^{(2)}_{V_k}(R_j\circ\Phi_{j,k},V_k)\|
 \le Cg_k L^{-s(k-j)}(g_k/g_j)^q\|R_j\|_{\rm graded}.
\tag{1}
\]
The norm includes quadratic/cubic kernels, anchor moments, quartic tail, and
original current/contact rows. The actual birth, physical and canonical radii,
promotion wait and once-owned rare payment remain attached to $R_j$.
LM42--47, pp.781--783, says that a previously integrated old mark must be carried
by its current representative, with all proper differences and filed projections.
Replacing it by a fine primitive would not preserve its saved age bound.

The primary source is T. Balaban, \emph{Large field renormalization II},
CMP 122 (1989), 355--392, (1.70)--(1.75), pp.~377--380.
Its component operation includes integrations and a history sum. On the large-field
components the old selectors and axial constraints remain. Removal of gauge fixing
on the other components is justified after the stated retained-field integration;
the paper distinguishes equivalence of densities from equality. Its positive
comparison and nonzero denominator concern the complete component operation.
These statements do not identify its conditional fluctuation law with the
single native reference law of \cite[Section 3]{ND}.

\paragraph{The exact connected coefficient and the missing row.}
Work at an actual real source-zero component where its original positive body
$Z_z=T_z1$ is nonzero. Normalize only this body, $P_z=Z_z^{-1}T_z$.
Let $\xi$ retain the history label and all older compact/constraint coordinates
in this component; let $B$ denote its last fluctuation coordinates. This notation
is a disintegration of the ORIGINAL measure, not replacement of its density.
Write $\pi_z(d\xi)$ and $\nu_{z,\xi}(dB)$ for its actual marginal and conditional
probabilities. For the admitted actual inherited observable $F$ and the distinguished
interaction $G$, with all their old readouts unchanged, finite-cutoff disintegration gives
\[
 \Cov_{P_z}(F,G)=
 \int\Cov_{\nu_{z,\xi}}(F,G)\,\pi_z(d\xi)
 +\Cov_{\pi_z}(m_F,m_G),\qquad
 m_F(\xi)=\nu_{z,\xi}F.                         \tag{2}
\]
In particular positivity bounds either variance summand by the full variance, not
conversely by the first summand. The second term is present even if $F$ is constant
in the last $B$ variable. Thus applying a fluctuation Poincare estimate, even if
valid for every fiber, cannot by itself estimate the full MT36 connected coefficient.

Equation (2) is an exact identified interface for the original operation. Neither
term has yet acquired the graded bound (1) on a general touched entry. In particular,
$\xi$ is not a subset of coordinates of ONE fixed convex product domain: it includes
actual discrete histories and older constrained variables. The depth-uniform
conditional-mean/current theorem in the native inherited gradient currents calculation therefore cannot be applied to the
second term merely by renaming $\xi$ as retained native coordinates.

\paragraph{Denominator and selector response, with the precise missing data.}
On any chart where the true conditional density can be written
$d\nu_{z,\xi}=Z_{z,\xi}^{-1}\chi_{z,\xi}(B)w_{z,\xi}(B)dB$,
the literal differentiated numerator/denominator identity is
\[
 D m_F=\nu(DF)+\Cov_\nu(F,D\log w)
 +Z_{z,\xi}^{-1}\int (F-m_F)w\,D\chi\,dB.       \tag{3}
\]
The last expression is a weak boundary distribution when the selector is sharp.
Any coordinate Jacobian, gauge determinant or reference response belongs to $w$
or to the original constraint distribution; it is not deleted. If constraints
cannot be eliminated on a common chart, (3) must instead be read as the derivative
of the original distributional numerator and denominator. For original outer
regular-pair differentiation at fixed integration coordinates, the source's
selectors are frozen; this does NOT assert they are constant in $B$ or in $\xi$.

The needed source-level bridge is consequently the explicit pre-gauge-removal
conditional form of R-II (1.70)/(1.71), using its R-I/IV selector definitions,
and the actual current representative in LM42 on that support. Native convexity
requires its actual domain and full log density, not just the denominator's
quadratic comparison. No verified equation presently equates this conditional
law to the sharp convex native reference. The source's integrated gauge equivalence
would suffice for a pushforward only after proving that every test of the chosen
observation belongs to its permitted gauge-invariant test class. A fixed outer
paired chart does not establish that condition by notation.

\paragraph{A genuine touched contact, and why it is only a diagnostic.}
The counterpart identifies the immediate stored physical contact (R12.8) of \cite[p.~1223]{W}:
\[
 b_{HH}=\frac{\sigma_j^2}{9\beta_j}
 \{\tfrac12(W_fP[h]+W_hP[f])-3(W_fD[h]+W_hD[f])
 +(2\beta_j)^{-1}(W_f\ell[h]+W_h\ell[f])\},
 \quad \sigma_j=g_j^{-2}a_j^{-2}.                 \tag{4}
\]
Choose overlapping completed stars actually touched by the original (1.60)
map $\Phi_B=\exp(i\eta H'_k(g_kB))U_k^0$. For its right-log tangent
$J_bv=(D_B\Phi_b[v])\Phi_b^{-1}$, the exact physical pullback derivative is
\[
 D_B(b_{HH}\circ\Phi)[v]
   =\sum_{b,A}(J_bv)^A(X_b^A b_{HH})(\Phi).       \tag{5}
\]
All original pointed test offsets and the Casimir action row remain in (4).
No native reference has been substituted in (5). Compact finite-word bounds
control the derivative with its original $\sigma_j^2/\beta_j$ factor, but neither
supply a bound for $J$ on every localized mixed wall nor restore a saved stable
age factor. This test checks a lower/contact row at its immediate substitution;
it is not the $R_j$ of (1) after earlier integration.

\paragraph{Current distinction.}
A Neumann current constructed inside $\nu_{z,\xi}$, if its Poincare hypothesis
were proved, represents only $F-m_F(\xi)$ on that fiber. It cannot represent
$m_F-\pi m_F$, since its fiber integral is zero. The second term in (2) thus needs
an actual inherited current/ordinary covariance estimate on $\pi$, retaining the
score and boundary terms (3). A pushforward of a native current also does not
supply an inverse tangent lift of the stored old physical current. Full weak
transport must retain the original kernel, selectors and live-reference derivatives.

\paragraph{Checks and status.}
The two terms in (2) follow by expanding $FG$ after adding and subtracting its
conditional means; no independence assumption is used. Formula (3) follows from
the quotient rule, so constants in $F$ cancel in its score and boundary terms.
The invariant derivative convention in (5) matches the right-log tangent of a
left multiplication path. These checks validate the interface, not the missing
bounds. No new MT36 norm clause is claimed closed. The next indispensable estimates are
actual selector/range admission for the same touched stored representative and a
bound for the retained-history term in (2), using one original reserve and the
saved age/radius data. Further native local moment refinements do not resolve it.
\begin{thebibliography}{9}
\bibitem{W} Wilson-source estimates, volumes 1--10 (\texttt{lattice\_theory/source\_estimates} in this collection), cited by source page.
\bibitem{ND} \emph{Nonlinear dispersion and unequal spatial weights}, \texttt{lattice\_theory/block\_and\_transfer/conditional\_estimates/nonlinear\_dispersion\_and\_unequal\_spatial\_weights.tex} in this collection.
\end{thebibliography}
\end{document}
